\chapter{考点方向四：自然语言处理与大模型理论 (Natural Language Understanding \& LLM)}
\label{ch:nlu}

\section{自然语言处理与大模型理论}

\source{本节对应考纲“统计建模、Word2vec、self-attention、LLM 预训练对齐、Scaling Law、基于 Embedding 的知识库推理”。主要参考 \textcite{jurafsky2026slp,goldberg2017nnnlp}；Transformer、Scaling Law、RAG、DPO 等补充参考 \textcite{vaswani2017attention,kaplan2020scaling,hoffmann2022training,lewis2020rag,rafailov2023dpo}。}

\subsection{统计语言模型}

\concepttarget{concept:ngram}

语言模型给序列 \(w_{1:n}\) 分配概率。链式法则给出

\[
  P(w_{1:n})
  =
  \prod_{i=1}^{n}P(w_i\mid w_{1:i-1}).
\]

\(n\)-gram 模型用 Markov 假设截断历史：

\[
  P(w_i\mid w_{1:i-1})
  \approx
  P(w_i\mid w_{i-n+1:i-1}).
\]

最大似然估计为

\[
  \widehat P(w_i\mid c)
  =
  \frac{\operatorname{count}(c,w_i)}
  {\operatorname{count}(c)}.
\]

问题是稀疏性：很多合理短语在训练集中没出现，MLE 会给零概率。因此需要 smoothing、backoff 或 interpolation。困惑度（perplexity）衡量测试序列平均预测难度：

\[
  \operatorname{PPL}(w_{1:n})
  =
  P(w_{1:n})^{-1/n}.
\]

PPL 越低，模型对测试文本赋予的平均概率越高。

\paragraph{\(n\)-gram MLE 的推导}

固定一个上下文 \(c\)。设词表为 \(\mathcal{V}\)，参数为

\[
  \theta_w=P(w\mid c),
  \qquad
  \sum_{w\in\mathcal{V}}\theta_w=1.
\]

给定计数 \(n(c,w)=\operatorname{count}(c,w)\)，该上下文相关的 log-likelihood 为

\[
  \sum_{w\in\mathcal{V}}
  n(c,w)\log\theta_w.
\]

引入拉格朗日乘子 \(\lambda\)：

\[
  \mathcal{J}(\theta,\lambda)
  =
  \sum_w n(c,w)\log\theta_w
  +
  \lambda\left(\sum_w\theta_w-1\right).
\]

一阶条件为

\[
  \frac{n(c,w)}{\theta_w}+\lambda=0,
  \qquad
  \theta_w=-\frac{n(c,w)}{\lambda}.
\]

对 \(w\) 求和并使用归一化约束：

\[
  1=\sum_w\theta_w
  =
  -\frac{1}{\lambda}
  \sum_w n(c,w)
  =
  -\frac{\operatorname{count}(c)}{\lambda}.
\]

所以 \(\lambda=-\operatorname{count}(c)\)，得到

\[
  \widehat P(w\mid c)
  =
  \frac{\operatorname{count}(c,w)}{\operatorname{count}(c)}.
\]

\subsection{Word2vec、GloVe 与 embedding}

\concepttarget{concept:word2vec}

Embedding 把离散词、实体或节点映射到向量空间：

\[
  e:\mathcal{V}\to\R^d.
\]

Word2vec skip-gram 用中心词预测上下文。给中心词 \(w\) 和上下文词 \(c\)，softmax 概率为

\[
  P(c\mid w)
  =
  \frac{\exp(u_c^\top v_w)}
  {\sum_{c'}\exp(u_{c'}^\top v_w)}.
\]

完整 softmax 太贵，negative sampling 用二分类目标近似：

\[
  \log\sigma(u_c^\top v_w)
  +
  \sum_{j=1}^{k}
  \E_{c_j\sim P_n}
  \log\sigma(-u_{c_j}^\top v_w).
\]

GloVe 从全局共现矩阵出发，使

\[
  u_i^\top v_j+b_i+\tilde b_j
  \approx
  \log X_{ij}.
\]

二者共同直觉是 distributional hypothesis：上下文相似的词语语义相似。

在 skip-gram 公式中，\(v_w\in\R^d\) 是中心词 \(w\) 的输入向量，\(u_c\in\R^d\) 是上下文词 \(c\) 的输出向量，\(P_n\) 是负采样分布，\(k\) 是每个正样本配多少个负样本。考试如果要求解释 negative sampling，应说明它把多分类 softmax 近似成“真实共现对 vs. 噪声共现对”的二分类问题，从而避免每次更新遍历整个词表 \(\mathcal{V}\)。

\subsection{Transformer self-attention}

\concepttarget{concept:self-attention}

Transformer 的输入是 token embedding 矩阵

\[
  X\in\R^{N\times d},
\]

其中 \(N\) 是上下文长度，\(d\) 是模型维度。单头 attention 先计算

\[
  Q=XW_Q,
  \qquad
  K=XW_K,
  \qquad
  V_{\mathrm{attn}}=XW_V.
\]

然后

\[
  \operatorname{Attention}(Q,K,V_{\mathrm{attn}})
  =
  \softmax
  \left(
    \frac{QK^\top}{\sqrt{d_k}}
  \right)V_{\mathrm{attn}}.
\]

每个 query 与所有 key 做相似度，softmax 得到权重，再对 value 加权平均。缩放因子 \(1/\sqrt{d_k}\) 防止点积方差随维度增大，使 softmax 过早饱和。

\paragraph{Attention 维度和缩放因子的推导}

设 \(W_Q,W_K\in\R^{d\times d_k}\)，\(W_V\in\R^{d\times d_v}\)。则

\[
  Q,K\in\R^{N\times d_k},
  \qquad
  V_{\mathrm{attn}}\in\R^{N\times d_v}.
\]

所以

\[
  QK^\top\in\R^{N\times N},
  \qquad
  \softmax(QK^\top/\sqrt{d_k})V_{\mathrm{attn}}\in\R^{N\times d_v}.
\]

若 query 和 key 的各坐标近似独立、均值为 \(0\)、方差为 \(1\)，则一个点积

\[
  q^\top k=\sum_{r=1}^{d_k}q_r k_r
\]

的方差为

\[
  \Var(q^\top k)
  =
  \sum_{r=1}^{d_k}\Var(q_r k_r)
  =
  d_k.
\]

因此未缩放的 logits 会随 \(d_k\) 变大，softmax 进入饱和区后梯度变小。除以 \(\sqrt{d_k}\) 后，点积方差回到常数量级，这是 scaled dot-product attention 的动机 \parencite{vaswani2017attention,jurafsky2026slp}。

自回归语言模型必须加 causal mask：

\[
  M_{ij}
  =
  \begin{cases}
    0, & j\le i,\\
    -\infty, & j>i.
  \end{cases}
\]

此时 attention 变成

\[
  \softmax
  \left(
    \frac{QK^\top}{\sqrt{d_k}}+M
  \right)V_{\mathrm{attn}},
\]

保证第 \(i\) 个位置不能偷看未来 token。这里 \(M\) 是加到 logits 上的 causal mask；padding mask 则遮蔽序列填充位置，二者不可混用。实现时还要核对标签是否已右移：若预测第 \(i\) 个目标 token，严格下三角掩码与输入/标签的移位约定必须匹配。多头 attention 把多个 head 的输出拼接后再线性投影。Transformer block 通常包括 attention、feed-forward network、residual connection 和 layer norm。

\subsection{LLM 预训练、指令微调与对齐}

\concepttarget{concept:pretraining-alignment}

自回归 LLM 的预训练目标是 next-token prediction：

\[
  \Lcal_{\mathrm{pretrain}}
  =
  -
  \sum_{i=1}^{n}
  \log p_\theta(w_i\mid w_{<i}).
\]

这就是交叉熵损失。它使模型学习语言统计、事实模式和许多任务的隐式格式，但目标本身只是预测文本，不保证有帮助、真实或安全。

这个目标同样来自最大似然。给定语料序列 \(w_{1:n}\)，自回归分解为

\[
  p_\theta(w_{1:n})
  =
  \prod_{i=1}^{n}p_\theta(w_i\mid w_{<i}).
\]

最大化 log-likelihood 等价于最小化

\[
  -\log p_\theta(w_{1:n})
  =
  -
  \sum_{i=1}^{n}
  \log p_\theta(w_i\mid w_{<i}).
\]

如果把真实下一个 token 写成 one-hot 分布 \(q_i\)，模型预测分布写成 \(p_i\)，单 token 损失就是交叉熵

\[
  H(q_i,p_i)
  =
  -
  \sum_{v\in\mathcal{V}}q_i(v)\log p_i(v)
  =
  -\log p_\theta(w_i\mid w_{<i}).
\]

指令微调（SFT）继续用交叉熵训练，但数据变成 instruction-response pairs：

\[
  \Lcal_{\mathrm{SFT}}
  =
  -
  \sum_{(x,y)}
  \sum_{t}
  \log p_\theta(y_t\mid x,y_{<t}).
\]

偏好对齐在上一节已给出 RLHF 与 DPO。它的局限包括 reward model 偏差、偏好数据覆盖不足、过度优化导致 reward hacking、长程事实性仍依赖检索或工具。

\subsection{Scaling Law}

\concepttarget{concept:scaling-law}

Scaling law 研究模型性能如何随参数量 \(N\)、数据 token 数 \(D\)、训练计算量 \(C\) 缩放。经验上，交叉熵损失常近似满足幂律：

\[
  L(N,D)
  \approx
  L_\infty
  +
  aN^{-\alpha}
  +
  bD^{-\beta}.
\]

这里 \(L_\infty\) 是不可约损失水平，\(a,b>0\) 是经验拟合常数，\(\alpha,\beta>0\) 是参数量和数据量对应的幂律指数。

这不是一个第一性原理定理，而是大规模实验规律。它的考试价值在于解释资源分配：固定 compute 时，参数量和数据量要平衡；只增大模型但数据不足会 data-limited，只增加数据但模型太小会 model-limited。Chinchilla 系列结果强调 compute-optimal 训练往往需要比早期大模型更多 token、更少参数。

\subsection{Embedding 知识库推理、RAG 与 MLN}

\concepttarget{concept:rag-mln}

知识库 embedding 把实体和关系映射到向量空间。TransE 的基本约束是

\[
  e_s+r\approx e_o
\]

对三元组 \((s,r,o)\) 成立。评分函数可写为

\[
  f(s,r,o)=-\|e_s+r-e_o\|.
\]

局限是多对多关系、组合逻辑、否定和量词难以表达。图神经网络、规则增强、路径推理和神经符号方法都是常见改进方向。

RAG 把检索器和生成器组合。给定查询 \(q\)，检索器从文档库 \(D\) 返回

\[
  R(q)=\{d_1,\ldots,d_k\},
\]

生成器根据 \(q\) 和检索文档生成答案：

\[
  p(x_{1:n}\mid q,R(q))
  =
  \prod_{i=1}^{n}
  p(x_i\mid q,R(q),x_{<i}).
\]

RAG 的价值是缓解幻觉、接入私有或时效知识、提供引用证据。主要风险是检索失败、检索噪声、上下文污染、重排错误和引用不忠实。工程闭环通常包括 chunking、embedding、向量索引、召回、重排、生成、引用和评估。

Markov Logic Network（MLN）用加权一阶逻辑公式定义 log-linear 分布：

\[
  P(X=x\mid E)
  =
  \frac{1}{Z(E)}
  \exp
  \left(
    \sum_j w_j n_j(x,E)
  \right).
\]

其中 \(n_j(x,E)\) 是第 \(j\) 个公式在世界 \(x\) 与证据 \(E\) 下满足的 groundings 数，\(w_j\) 是该公式的权重，\(Z(E)=\sum_x \exp(\sum_j w_j n_j(x,E))\) 是在证据 \(E\) 下的配分函数，用来保证概率和为 \(1\)。权重大表示强偏好；无限大权重可视为 hard constraint。它和 embedding 方法的差别是：MLN 显式编码逻辑约束，embedding 方法更擅长连续相似性和大规模近似。

\paragraph{考核内容}

\begin{itemize}
  \item 从链式法则写出语言模型目标，解释 \(n\)-gram 的 Markov 假设与 smoothing。
  \item 推 Word2vec skip-gram 或 negative sampling 目标，解释 embedding 的语义来源。
  \item 写出 self-attention 的 \(Q,K,V\) 公式、维度、mask 和复杂度。
  \item 解释 next-token pretraining、SFT、RLHF、DPO 的目标差异。
  \item 说明 scaling law 的经验形式和 compute-optimal trade-off。
  \item 比较知识库 embedding、RAG 和 MLN 的能力边界。
\end{itemize}

